\section{Award-Layer Triage}
\label{sec:award_screen}

The award layer observes outcomes before it observes conduct. A
losing bid may be sincere, poorly priced, capacity constrained,
exploratory, strategic, or coordinated. The screen therefore
uses what the award layer can measure at the triage stage: which
firms appeared, which firms won, and which firms kept appearing
without ever winning. The object is a ranking for forensic
priority, not proof of conduct or a classification of cartel
membership. This is the law-enforcement-economics decomposition
of investigative attention into ranking and proof
\citep{becker1968crime,stigler1970optimum}: when proof is costly,
the question is where to begin, not what to assume.

This section defines that ranking and the empirical implications
that discipline it. Subsection~\ref{sec:award_primitive}
introduces persistent zero-win participation as the award-layer
primitive. Subsection~\ref{sec:award_rule} explains why the
continuous ranking can be converted into an operational queue or
binary frequent-loser flag. Subsection
\ref{sec:testable_implications} states the tests that the rest
of the paper uses to evaluate whether the ranking is informative
for enforcement.

\subsection{Persistent Zero-Win Participation}
\label{sec:award_primitive}

Let $T_i$ denote the number of BEC tender-items in which firm
$i$ participates, and let $W_i$ denote the number it wins. The
award-layer screen operates inside the zero-win stratum,
$\{i:W_i=0\}$, which contains $\valAlwaysLosers$ firms. Within
that stratum, the continuous score is
\[
  s_i=\log(1+T_i).
\]
The log transformation stabilizes the scale but does not change
the ordering. The economic object is the rank: among firms the
award layer observes only as losers, which ones appear
persistently enough to warrant bid-layer follow-up?

The maintained condition is monotonic rather than classificatory.
If cartel-deployed losing participants are present, they should
appear more persistently in the zero-win stratum than ordinary
losing firms, because such a role is useful only through repeated
deployment in cartel-targeted tenders. The condition is
maintained, not derived. \ref{app:framework_submission}
formalizes it as a monotone-likelihood-ratio ranking result in
the tradition of \citet{karlin1956mlr}. The
main text uses it more modestly: persistent zero-win participation
is a cheap statistic that can rank loser-side exposure under a
transparent behavioral assumption and direct attention to the bid
layer where conduct can be evaluated. It does not assign a firm
type or cartel membership.

The zero-win restriction is therefore a scope restriction. Many
firms never win for ordinary competitive reasons---poor fit, high
costs, capacity constraints, market exploration, or simple bad
luck---and the screen's content comes from persistence inside that
heterogeneous loser-side set. If the ranking has enforcement
value, it should concentrate adjudication-anchored loser-side
adjacency beyond what participation volume, product-market
overlap, and opportunity to meet CADE anchors mechanically
produce.

\subsection{From Ranking to an Operational Queue}
\label{sec:award_rule}

The continuous score is the central object because agencies can
translate a rank into different investigative protocols. A team
with fixed capacity could inspect the top $k$ firms. A
monitoring unit could refresh the rank each quarter and examine
bunching below old cutoffs. A case team could use the score to
request LANCES bid files for a survivor pool before applying
bid-distribution forensics. These implementations differ, but
all preserve the same sequencing logic: award records first,
costly bid recovery second.

For an auditable binary implementation, we also report a
frequent-loser flag. The rule marks zero-win firms whose
participation count exceeds the median plus $1.5$ times the
interquartile range of the always-loser participation
distribution. In BEC, the statistical cutoff is
$\valThresholdStat$, implemented as $T_i\geq\valThreshold$, which
flags $\valFL$ firms, or $\valFLrateAL$ of the always-loser
stratum. The cutoff is fixed from the participation distribution
alone, before any CADE label enters the analysis.

The cutoff is an auditable administrative implementation of the
rank, not a structural or legal threshold. The paper does not
rest on the number $14$; it rests on the continuous loser-side
participation ranking. Two analysts applying the same rule to
the same award records obtain the same first-stage list. Above
the cutoff does not mean cartel member; below it does not mean
clean. Because a static threshold can be gamed, the enforcement
architecture treats the binary flag as one implementation among
several: agencies can also use continuous ranks, top-$k$ queues,
refreshed thresholds, or budget-calibrated survivor pools, in
line with the operational-screening design framework articulated
in \citet{sanchezgraells2019screening} and the robust-screen
principles of \citet{chassang2022robust}. The
strategic-adaptation analysis later returns to this point.

\subsection{Testable Implications}
\label{sec:testable_implications}

The screen is useful only if it survives the threats implied by
its limited information. We follow the threat-driven screen-design
tradition of \citet{harrington2008detecting} and the
bid-distribution screen evaluation in \citet{imhof2018screening}:
a screen is informative only to the extent that it survives the
alternative explanations a thoughtful referee would offer.
Table~\ref{tab:testable_implications_submission} organizes the
validation map into six implications---loser-side ranking, scope
asymmetry, economic profile, complementarity, gatekeeping, and
strategic response---each paired with its prediction, test, and
main threat. The core claim is that persistent zero-win
participation concentrates adjudication-anchored cobidders within
the always-loser stratum, against the principal threat of
mechanical exposure (participation volume and the opportunity to
meet CADE anchors); the remaining rows discipline the loser-side
scope asymmetry against direct defendants, the economic content of
the cobidder profile, the award--bid complementarity, the
gatekeeping cost-of-evidence, and fragility to strategic
adaptation. None of these predictions defines legal membership or
proof.

\begin{table}[htbp]
\centering
\caption{Testable implications and validation threats}
\label{tab:testable_implications_submission}
\footnotesize
\begin{adjustbox}{max width=\textwidth,center}
\begin{threeparttable}
\begin{tabular}{p{3.2cm}p{4.4cm}p{4.3cm}p{3.2cm}}
\toprule
Implication & Prediction & Test & Main threat \\
\midrule
Loser-side ranking & Persistent zero-win participation ranks always-loser cobidders. & Conservative benchmark, timing audit, exposure-adjusted validation. & Participation volume, market overlap, and opportunity to meet CADE anchors. \\
Scope asymmetry & Performance is weaker for direct defendants than for cobidders. & Direct-defendant AUC and temporal direct-defendant check. & Misreading the screen as generic cartel-membership detection. \\
Economic profile & Cobidders differ from other frequent losers in loser-side exposure and bid patterns. & Firm-level and bid-level profile comparisons. & Product-market composition and ordinary high-volume losing. \\
Complementarity & Award-layer signal adds information to bid-layer forensics. & Same-target award, bid, and combined AUCs. & Full-observability benchmark confused with operational sequencing. \\
Gatekeeping & Award-first sequencing reduces bid-microdata recovery with limited recall loss. & Sequential survivor-pool exercise and cost-recall frontier. & Survivor-pool choice and arbitrary budget choice. \\
Strategic response & Static cutoffs are more fragile than ranks or refreshed queues. & Adaptation diagnostics. & Simulation is not an equilibrium model. \\
\bottomrule
\end{tabular}
\begin{tablenotes}
\footnotesize
\item \textit{Notes:} Predictions evaluate whether the award-layer
ranking is useful for evidence allocation. They do not define
legal membership or proof.
\end{tablenotes}
\end{threeparttable}
\end{adjustbox}
\end{table}

Together, these tests evaluate the screen at the stage for which
it is designed: the allocation of costly forensic attention
before liability evidence is assembled. The BEC cutoff is an
implementation detail; the portable object is the staged
architecture that uses cheap administrative records to discipline
where costly cartel forensics should begin.
